\section{\large 高校数学でのベクトル}
高校数学のベクトルと大学数学のベクトルに違いがあるのかというわれると,ないと感じている.高校数学では「平面ベクトル」「空間ベクトル」までを扱った.
これはつまり,3次元空間までのベクトルということである.それは図で書くことができる最大の次元であるからである.イメージしやすいのがメリットだったが大学数学ではn次元のベクトルまで扱う.しかし,それだけのように感じる.
私が線形代数学を学習するとき,定義はn次元を想定して描かれているが3次元の図を基本イメージして理解に努めていた.4次元ってどうイメージするのか？私にはわからない.\\

まず高校数学でのベクトルでは$\vec{a}=(1,2,3)$と書いていただろう.ただ,大学では縦書きになる.そして,ベクトルを表していた文字の上にあった矢印はなくなり,代わりに\textbf{太く}書く.\\
つまり,
\[
\vec{a} = (1,2,3) \qquad \rightarrow \qquad \symbfit{a} = \begin{bmatrix}
    1 \\2\\ 3 \end{bmatrix} \text{もしくは} \begin{pmatrix}
        1 \\2\\ 3 \end{pmatrix}\text{となる.}
\]
好みだと思うが,私は[\quad]を使う.注意点は$\begin{vmatrix}1 \\ 2 \\ 3\end{vmatrix}$は間違いですので気をつけるように.これは「行列式」で使う書き方で,意味が変わってしまう.
行列式はSection.\ref{sec:det}で学ぶ.
\section{線形結合}
この言葉は高校数学ではなかった.実はこれは高校生の頃に何気なくやっていることに名称がついただけの話である.\\
\underline{線形結合は「k倍したベクトル同士を足せますよ」.これだけの話.}
\[
c\symbfit{v}+d\symbfit{w}=c\begin{bmatrix}1 \\ 1\end{bmatrix}+d\begin{bmatrix}2 \\ 3\end{bmatrix} =
\begin{bmatrix}c+2d \\ c+3d\end{bmatrix}
\]
この書き方は初学者には特別な変形に見えるかもしれないが,今まであなたがやっていたように\\
ベクトルを横書きにして( \quad )に書き直して同じ計算をしてみるといい.表記の違いだけということに気づくだろう.\\
そして上から数字を第1要素,第2要素...なんて読んでいく.これはx軸方向に,y軸方向にとか言っているとz軸方向で詰んでしまうからだと思ってほしい.\\
そうなると第2要素までのベクトルを2次元ベクトル,第n要素まであるベクトルはn次元ベクトルとなるね.\\
事前に言っておくが,4次元以降のベクトルは想像できないので,割り切って紙面上数字と思って学習した方が脳に優しいのでおすすめ.\\
\begin{tcolorbox}[colframe=black!80,colback=white,title=線形結合$_{\text{linear combination}}$]
$c\symbfit{v}とd\symbfit{w}の和は,\symbfit{v}と\symbfit{w}の線形結合である.$    
\end{tcolorbox}
必要な道具はここで揃った.実は新しい内容である「\textbf{行列}」に入るにはもうこれで十分である.意外と身構えることではない.
\section{(重要)行列}
\subsection{行列への拡張}
行列の感覚としては最も効率よくベクトルを並べる方法とでも思っている.確かに書き方は変わるが「線形結合」を意識すると難易度はグッと下がる.どんなものか楽しみにしておくとよい.\\
行列へと自然に持っていくために少し脱線する.\\
皆さんはこの問題は解けるだろうか？\\
\[
\begin{array}{ccccccc}
    x & + &2y & +&3z &= &6 \\
   2x & +&5y & +&2z &= &4 \\
   6x & - &3y & +&z &= &2 \\
\end{array}
\]
おそらく,「3つの式で成り立つx,y,zを求める問題」から「3つの面が交わってできた1点を求める問題」の範囲に収まる認識だと思う.
しかし,線形代数学を学習した学生はそうは考えない.考えてはいけない.なぜならベクトルで考える方がはるかにイメージしやすいからである.
上の式をx,y,zで括ろうと思えばこのように括ることができる.\\
\begin{equation}\label{eq:linear_comb}
x
\begin{bmatrix} 1 \\ 2 \\ 6 \end{bmatrix}
+y
\begin{bmatrix} 2 \\ 5 \\ -3 \end{bmatrix}
+z
\begin{bmatrix} 3 \\ 2 \\ 1 \end{bmatrix}
=
\begin{bmatrix} 6 \\ 4 \\ 2 \end{bmatrix}
\end{equation}
うだうだ書こうとは思ったが,共通因数であるx,y,zを縦に括るともう[ \quad ]で括るほうが見やすく,
そうするともうベクトルの世界に式が入ってしまった.\\
これによって式は,
$\begin{bmatrix} 1 \\ 2 \\ 6 \end{bmatrix}$をx倍,
$\begin{bmatrix} 2 \\ 5 \\ -3 \end{bmatrix}$をy倍,
$\begin{bmatrix} 3 \\ 2 \\ 1 \end{bmatrix}$をz倍したベクトルたちの和が作るベクトルが
$\begin{bmatrix} 6 \\ 4 \\ 2 \end{bmatrix}$となるようにx,y,zを求めたいという式となる.
\begin{figure}[h]
    \centering
    \includegraphics[width=0.6\textwidth]{cal_vec.png}
    \caption{連立方程式のベクトル空間的解釈}
    \label{fig:sample}
\end{figure}\\
※「列ベクトル」,「第k列」はここでは無視で\\
どうだろう？立体空間に傾斜もわからない3つの平面を交わらせて1点をイメージするより,
3本の矢印を伸ばして,その和を右辺のベクトルに持っていく作業と考える方がイメージしやすいとは思わないだろうか？\\
図\ref{fig:sample}にもあるが答えは
$\begin{bmatrix} x \\ y \\ z \end{bmatrix} = 
\begin{bmatrix} 0 \\ 0 \\ 2 \end{bmatrix}$である.
実際に式(\ref{eq:linear_comb})に代入して確かめてみよう.
\[
0
\begin{bmatrix} 1 \\ 2 \\ 6 \end{bmatrix}
+0
\begin{bmatrix} 2 \\ 5 \\ -3 \end{bmatrix}
+2
\begin{bmatrix} 3 \\ 2 \\ 1 \end{bmatrix}
=
\begin{bmatrix} 6 \\ 4 \\ 2 \end{bmatrix}
\]
ん？これ線形結合とは思いませんか？\\
ぬるっと大学数学の「線形結合」と合流する.(そう見えたら狙い通り)\\
ここに来て頑張りどころが初めてくる.ここで行列の書き方と行列の積を紹介する.ここは定義として飲み込んでほしい.
行列同士の計算方法はSection.\ref{sec:matrixRule}で説明するため安心してほしい.
\begin{tcolorbox}[colframe=black!80,colback=white,title=行列とベクトルの積]
\[
\begin{bmatrix}
    & & \\
    \:\symbfit{u}& \symbfit{v} & \symbfit{w} \: \\
    & & \\
\end{bmatrix}
\begin{bmatrix}c \\ d \\ e \end{bmatrix}
=
c\symbfit{u}+d\symbfit{v}+e\symbfit{w}
\]
左の[ \quad ]にあるのはベクトルである.今回の問題でにおいては,
\begin{equation}\label{eq:lenar_equal}
0
\begin{bmatrix} 1 \\ 2 \\ 6 \end{bmatrix}
+0
\begin{bmatrix} 2 \\ 5 \\ -3 \end{bmatrix}
+2
\begin{bmatrix} 3 \\ 2 \\ 1 \end{bmatrix}
=
\begin{bmatrix} 6 \\ 4 \\ 2 \end{bmatrix}\quad\Longleftrightarrow\quad
\begin{bmatrix}
    1 & 2 & 3 \\
    2 & 5 & 2 \\
    6 & -3 & 1
\end{bmatrix}
\begin{bmatrix}
    0 \\ 0 \\ 2
\end{bmatrix}
=
\begin{bmatrix}
    6 \\ 4 \\ 2
\end{bmatrix}
\end{equation}
\end{tcolorbox}
上の定義とじっくり照らし合わせてほしい.各ベクトルが縦書きを保ったままぎゅっと1つの[ \quad ]にまとまった形をしている.
これを一般的に「\textbf{行列}(matrix)」と呼ばれる.イメージは「行を成す数字の羅列」を略して「行列」であろう.線形結合から行列に書き直すのは慣れればどうってことないが,
初学者にはこの独特な掛け合わせが行列を難しく感じさせる.最初の見慣れた連立方程式から線形結合を経て行列への数字の動きを書いてみるといい.\\
私も最初は難しく感じた.というのも線形結合を飛ばして行列をみると3\times 3の数字が何を表しているのかすぐにはわからないからだ.
\underline{行列から線形結合.線形結合から行列を繰り返し考えたり書いたりして慣れてみ}\\\underline{よう.}

さて,図\ref{fig:sample}が上では挿入されており,「列ベクトル」とある.これについて説明する.教科書$^{\cite{strang}}$では先に列ベクトルを説明するため先に出ているが,
私個人としては行列を先に見てからの方が納得できると思ったからである.
行列は横方向に「行」,縦方向に「列」と呼ぶ.
例えば第1行は(1,2,3)であり,第1列は(1,2,6)である.日本人限定の恩恵だが「行」という字の二の部分が方向と一致する.
「列」も同様である.有効な覚え方である.式(\ref{eq:lenar_equal})においては,(-3)は3行2列の要素なんて言ったりする.改めて図\ref{fig:sample}を見てみよう.これらのベクトルは後に行列の列となる.
そこから「列ベクトル」と名称がつくとなると,なぜ「列ベクトル」なんだという疑問はなくなるのではないだろうか？
今回の3行3列からなる行列を「$3\times3$の行列」と言ったりする.行列が何行何列かは行列の「\hypertarget{size}{\textbf{サイズ}}」という.ちょっとしたことを言うと,列ベクトルも3行1列の行列とも言える.だが,列ベクトル呼びで結構.
\\

これから行列主体の内容になっていくが,\textbf{わからなくなったら,左辺の行列と列ベクトルの積は線形結合であることを忘れないでほしい.}
大学の講義における線形代数学のゴールは「対角化(Section.\ref{sec:diagonalization})」であるが,これ本質的に理解させようとすると,この行列の積を線形結合であることに慣れておくことが大事だと考えているからだ.
このテキストでは暗記で終わらせるつもりはない.また,グラムシュミット(Section.\ref{sec:gram})に限らず,これからの行列の計算では「\textbf{わからなくなったら線形結合まで帰着する.}」ことで解決することが多いと信じている.
私自身の現在のレベルまでの学習において,これで乗り越えてきている.
そのため,「列ベクトル」から行列の解説を始めることには\underline{大きな価値があると考えている.}
\subsection{転置(transpose)}
転置とは\textbf{行と列の位置関係を交換すること}である.本によって行列に対して「転置」をしたことを示す書き方が異なる.
\begin{enumerate}
    \item ${}^t\!A$\;(texでは\verb|{}^t\!A|と書く)
    \item $A^t$(小文字)
    \item $A^T$(大文字)
\end{enumerate}
本テキストでは$A^T$と表現している.

例えば,$3\times5$の行列$A$を転置した$A^T$は$5\times3$となる.
\[
A=\begin{bmatrix}
    11 & 12 & 13 & 14 & 15 \\
    21 & 22 & 23 & 24 & 25 \\ 
    31 & 32 & 33 & 34 & 35 
\end{bmatrix}\qquad\to\qquad
A^T=
\begin{bmatrix}
    11 & 21 & 31 \\
    12 & 22 & 32 \\
    13 & 23 & 33 \\
    14 & 24 & 34 \\
    15 & 25 & 35 
\end{bmatrix}
\]
また,列ベクトルに対してももちろん転置は可能である.
\[
    \symbfit{a}=\begin{bmatrix}
        1 \\ 2 \\ 3
    \end{bmatrix}\qquad\to\qquad
    \symbfit{a}^T=\begin{bmatrix}
        1 & 2 & 3
    \end{bmatrix}
\]
\bigskip\\
転置に対しては一般的に覚えておくべき重要な法則が存在するため,紹介する.しかしそれには\underline{行列同士の計算}\\\underline{の知識が必要}である.セクションが前後してしまって申し訳ないが,初学者や行列同士の計算覚えてない,読んでもピンとこなかった方は先にSection.\ref{sec:matrixRule}を先に読むことを\textbf{強く勧める}.
\begin{itemize}
    \item 行列同士の足し算の結果を転置するのと,転置してから足し算しても足される数同士の組み合わせは変わらないので,
\[
    (A+B)^T=A^T+B^T
\]
\item 和と同様,積を転置するのと,転置して積をとっても,組み合わせは変わらない.しかし,転置すると行列のサイズが変わり,積が取れなくなることがある.\\
例を挙げる.$3\times4$と$4\times2$の積において転置すると,$4\times3$と$2\times4$でたしかに積が取れない.そこで,場所を入れ替えると$2\times4$と$4\times3$なって積が取れる.
\[
    (AB)^T = B^TA^T
\]
\end{itemize}
実際に簡単な行列を作って検証してみると一層理解が深まるだろう.また,例では2つの行列だったが,n個の行列でも成り立つ.
\begin{tcolorbox}[colframe=black!80,colback=white,title=転置の性質]
\begin{align*}
    (A_1+A_2+\dots+A_n)^T&=A_1^T+A_2^T+\dots+A_n^T\\
        (A_1A_2\dots A_n)^T&=A_n^T\dots A_2^TA_1^T
\end{align*}
\end{tcolorbox}
\section{線形方程式の求解}
\subsection{(重要)簡約化}
そういえばベクトルは$\symbfit{a}$と書くなら行列にも何か代表する文字があってもいいはず.あります.小文字のアルファベットが列ベクトルならそれからなる\underline{行列はAとなる.}特に決まりがない場合はAが使われる.\\
先ほどの問題で考える,シンプルに
\begin{equation}\label{eq:Ax_b}
\begin{bmatrix}
    1 & 2 & 3 \\
    2 & 5 & 2 \\
    6 & -3 & 1
\end{bmatrix}
\begin{bmatrix}
    x \\ y \\ z
\end{bmatrix}
=
\begin{bmatrix}
    6 \\ 4 \\ 2
\end{bmatrix}
\qquad \leftrightarrow \qquad A\symbfit{x} = \symbfit{b}
\end{equation}
と書ける.Aは太字じゃないので注意.

さて,私たちはこれから何をするのだろうか？それは$\symbfit{x}$を求める方法を学ぶことである.\\
先ほどの問題ではすぐ答えを言っており,解き方の紹介はしていない.答えが予想しやすいので,違和感を感じる人はいなかったかもしれないが,
連立方程式ではこの$\symbfit{x}$にあたる数を求めるのが大変であった.せっかく行列表記にしたのだから,それを活かした新しい画期的な解き方でないとおかしい.\\
答えを言うと「微妙」である.私はこれから「\textbf{拡大係数行列}」から始まり,「\textbf{簡約化}」へ触れていくことになるが,計算が辛い.計算ミスもする.私は計算ミスがとても多くこの単元では「良」だった.
計算ミスを減らす方法はあるかと言われると,見つかっていない.しかし,この内容が後の内容につながっていく.具体的には「\textbf{逆行列}」「\textbf{行列式}」「\textbf{余因子行列}」である.
なので,苦行は今後のためのツールだとして頑張って欲しい.\\
※内情を言うと,行列にすると「\textbf{前進消去と後退代入}」という画期的な解き方へアクセスできるがテストに出ない.ただ,知っていると早く解けるのでSectin.\ref{sub:Gaussian}(発展)で紹介する.そこでは「消去の破綻」にも触れる.\\
\\

まず用語から.Aは連立方程式の係数からなるので「\textbf{係数行列}」と特に言われる.「拡大係数行列」は字の如く,係数行列を拡大した書き方である.
具体的に行列(\ref{eq:Ax_b})で例を示すと,[A $\symbfit{b}$]こんな感じ.今回の問題では\\
\[
\begin{bmatrix}
    1 & 2 & 3 & 6 \\
    2 & 5 & 2 & 4 \\
    6 & -3 & 1 & 2
\end{bmatrix}
\]
つまり行列$A$の要素に$\symbfit{b}$の要素が横にくる行列を拡大係数行列という(「拡大行列」ともいう).なぜそのような書き方をするかであるが,
これから私たちは行列に対して「簡約化」と言うものを行うことで$\symbfit{x}$を求める.拡大係数行列で書いて処理した方が見やすいからである.
\bigskip\\
「簡約化」は先にゴールを示した方がわかりやすいと思う.簡約化という「処理」を行うことで,どうなるかを示す.
上記にある拡大係数行列を使っても良いのだが,数の並びの都合上見づらいので,新しく拡大係数行列を用意した.簡約化もそれを使って説明する.
\begin{framed}
\[A=
\begin{bmatrix}
    2& 4 & -2 & 2 \\
    4 & 9 & -3 & 8 \\
    -2 & -3 & 7 & 10
\end{bmatrix}
\quad\rightarrow\quad R=
\begin{bmatrix}
    1& 0 & 0 & -1 \\
    0 & 1 & 0 & 2 \\
    0 & 0 & 1 & 2
\end{bmatrix}
\]
\begin{center}
    　　簡約化\\
\end{center}
\vspace{-4mm}
\end{framed}
Aの対角成分が1でそれ以外が0となった.これが簡約化である.ここまでくるとメリットが見えてくる.簡約化された行列は\underline{大変解きやすいのだ.}
例えばこのような連立方程式を解くことを考えてみる.
\[
\begin{array}{ccccccc}
    1x & + &0y & +&0z &= &-1 \\
   0x & +&1y & +&0z &= &2 \\
   0x & + &0y & +&1z &= &2 \\
\end{array}
\]
全ての式が1つの変数に対する1つの値から式となっており,答えはすぐに$x=-1,y=2,z=2$と求まる.\\
この連立方程式を行列で表すと,
\begin{equation}\label{eq:solableByRREF}
\begin{bmatrix}
    1& 0 & 0 \\
    0 & 1 & 0 \\
    0 & 0 & 1 
\end{bmatrix}
\begin{bmatrix}
    x \\ y \\ z
\end{bmatrix}
=
\begin{bmatrix}
    −1 \\ 2 \\ 2
\end{bmatrix}
\end{equation}
さらに拡大係数行列にすると,行列$A$の簡約化後である行列Rと一致する.これによって簡約化によって解きやすくなったことがわかったと思う.

それではその肝心な簡約化の方法であるが,実は皆さんが連立方程式を解くときとそこまで変わらない.中学生でやったようなxとyからなる連立方程式では解き方が2種類あったはず.
1つは$x=,y=$と式変形をして代入する方法.もう1つは「上の式を2倍して,下の式と引く」といった操作を行なって文字を減らす方法.\textbf{簡約化は後者をやっているに過ぎない.}\\
\\
具体的な流れを示す.
簡約化後の拡大係数行列をみると,Aにあたる部分の対角成分が1となっている.この位置の要素をここではpivot(ピボット)と呼ぶ.pivotの定義はこの後出る.体格要素をすべてpivotと呼ぶわけではないので注意.
基本的にはpivotの位置にある要素を1にして,それ以外の要素を0にしていく工程が続く.
まず1行1列は2であるので,第1行を2でわる.
\[
\begin{bmatrix}
    1& 2 & -1 & 1 \\
    4 & 9 & -3 & 8 \\
    -2 & -3 & 7 & 10
\end{bmatrix}
\]
これでpivotは1になった.もう1度簡約化後の拡大係数行列をみると,1行1列の1より下の要素は0なので,0を作る.
第1行を-4倍したものを第2行に足す.(第1行を4倍したものを第2行に引くでも同じ).第1行を2倍したものを第3行に足す.
\[
\begin{bmatrix}
    1& 2 & -1 & 1 \\
    0 & 1 & 1 & 4 \\
    0 & 1 & 5 & 12
\end{bmatrix}
\]
これで第1列の処理は終了\\
次に第2列の処理をする.2行2列の要素を1にするのだが,今回はたまたまpivotが1になったくれたので上下の要素を0にする.
やり方は「式を定数倍して足す」という先程と同じことをする.ポイントは第1列には0があり,いくら式同士でかけたり足したりしても第1列の値は変わらないということ.
つまり,第1列を気にすることなく,第2列の簡約化に集中することができる.
具体的な処理は第2行を2倍して第1行に足す.第2行を−1倍して第3行に足す.
\[
\begin{bmatrix}
    1& 0 & -3 & -7 \\
    0 & 1 & 1 & 4 \\
    0 & 0 & 4 & 8
\end{bmatrix}
\]
第3列も同様に処理すると
\[
\begin{bmatrix}
    1& 0 & 0 & -1 \\
    0 & 1 & 0 & 2 \\
    0 & 0 & 1 & 2
\end{bmatrix}
\qquad \text{\textcolor{red}{簡約化完了}}
\]
ここで注意点をいくつか
\begin{itemize}
    \item 大学の講義では書き方が違うがそちらを推奨する.これは解説用に書き方を変えている.
    \item 今回の簡約化は整数のみとなるが,途中や簡約化後の第4列に分数が現れる方が多い.
    \item 初心者向けに対角成分を1にすることから始めているが,慣れればそうしない方が早い時があることに気づく.
    今回の問題でいうなら,第1行を最初に2で割る前に,第1行を−2倍して第2行へ足す.第1行を1倍して第3行に足す.\\
    これを先に行った方が素早く0が現れるし,計算ミスが少なくなる.
\end{itemize}
他にも行を入れ替えるなど手段がある.例えば第1行の数字と第2行の数字を交換する.行の入れ替えをしても大丈夫なのか？と思った人もいるかもしれないが,
連立方程式の時を思い出してほしい.式の順番を変えても解には影響ない.連立方程式の形で式を交換する.それを行列にする.これは行の入れ替えをしたという.小技はたくさんある.
つまり,簡約化にはセンスや慣れによって速度や過程の形に差が生まれる.講義や先生が出す解説動画では行列の形で解き方が変わっていくので,
それを見て小技を蓄えていくといいだろう.
\newpage
\subsection{(重要)rank,階数なるもの}
rankは簡約化の問題とセットで出題されている.本当は「列空間」(Section.\ref{subsub:colspace}(発展))というものを学んでからの方がrankが何を指しているのかわかりやすい.発展内容であるため,ここでは「何故」は飛ばして結論だけ言っておしまいにする.
\begin{tcolorbox}[colframe=black!80,colback=white,title=rank]
\begin{center}
    \textbf{pivotの個数}
\end{center}
\begin{framed}
簡約化後の各行で左から右へ見ると,最初に現れる非零が存在する行がある.その個数.\\\\
注意点としては拡大係数行列でのみrankが存在するわけではない.行列が単体でポンと出題され,「rankを求めよ.」と言われることは十分ある.
簡約化して数えればいい.それで求まる.
\\
今回ならrankは3である.いくつか例を挙げる.
\[
A=
\begin{bmatrix}
    1 & 0 & 0 & 2 \\
    0 & 1 & 0 & -2/5
\end{bmatrix} \qquad
\therefore rank(A) = 2
\]
\[
B=
\begin{bmatrix}
    1 & 0 & 0 & 1 \\
    0 & 1 & 0 & 5 \\
    0 & 0 & 0 & 0
\end{bmatrix} \qquad
\therefore rank(B) = 2
\]
rankの説明で
「簡約化後の各行で左から右へ見ると,最初に現れる非零が存在する行がある.」とあるが,この非零の数をpivot(ピボット)という.非零と言っているが簡約化しているので,ほぼ1のことを言っている.情科の学生なら必ずいつか聞くことになる言葉.
\end{framed}
\end{tcolorbox}
pivotのより厳密な説明は(発展)で触れるがおすすめ.

\subsection{(発展)ガウスの消去法と列空間}
\subsubsection{行列であることのメリット(前進消去と後退代入)}\label{sub:Gaussian}
テストでは「簡約化せよ」と出題されるので,簡約化されてから,式(\ref{eq:solableByRREF})が示すように線型方程式の解が求められる.加えて,$rank$も求められた.しかし,簡約化は手間がかかるし,
簡約化しなくても求められたら,簡約化したくないものだ.簡約化せずとも求まるのは,
\begin{itemize}
    \item 線形方程式の解$\symbfit{x}$
    \item rank
\end{itemize}
である.つまり,\textbf{簡約化は不要なのだ！}
\paragraph{線形方程式の解}
「\textbf{前進消去と後退代入}」の話をしよう.簡約化と手順は大きく変わらないが,「処理」の回数が少なければ計算ミスも減るだろう.\\
簡約化されたpivotが1でそれ以外が0になった行列をRREF行列(Reduced Row Echelon Form)と言う.\\名前がカッコいいからと言う理由で紹介する.

この\underline{前進消去と後退代入は二部構成となっている.}前進消去を行い,後退代入をすることで解が出る.
\newpage
まず,\textbf{前進消去}
\[
\begin{bmatrix}
    2& 4 & -2 & 2 \\
    4 & 9 & -3 & 8 \\
    -2 & -3 & 7 & 10
\end{bmatrix}
\rightarrow
\begin{bmatrix}
    2& 4 & -2 & -2 \\
    0 & 1 & 1 & 4 \\
    0 & 0 & 4 & 8
\end{bmatrix}
\]
\begin{center}
    　　前進消去
\end{center}
簡約化と同様にゴールを先に示した.簡約化より非零の要素が多く残る形であるが十分である.簡約化との違いは2つ
\begin{enumerate}
    \item \textbf{pivotを1にしなくていい}
    \item \textbf{0にしなくてはいけないのはpivotの下側のみ}
\end{enumerate}
ポイントなのはこの形にするまでの工程は簡約化と同じ工程をするが,pivotを1にする必要がないし,行同士の計算が少ない.注目して欲しいのはpivotが
階段のようにギザギザしている.綺麗な階段にならずとも,行列を斜めに切り,下三角形が全て0になるようにすれば,階段になる.これを「階段化」と言う.
例えば,
\[
\begin{bmatrix}
    1 & 2 & 0 & 3 & 4 \\
    0 & 0 & 1 & 4 & 5 \\
    0 & 0 & 0 & 2 & 1 \\
\end{bmatrix}
\]
これも階段化に成功している.また,これは前進消去が完了した形とも言える.このような行列を「行階段行列」と言います.名前の由来がわかってくれたかな？\\
※「行階段行列」は英語で「REF行列(Row Echelon Form)」という.名前がカッコいいk\\
また,消去の過程で次操作をしたい行が全て0になってしまったらどうしよう？これはそれより下から都合の良い行を見つけて交換してあげればいい.
都合の良い行とはpivotが存在する行のことである.\\
「階段」という言葉は同じ
階段化は簡約化の行列でも言える.RREF行列は日本語では「簡約階段行列」という.
個人的にはREF行列とRREF行列の違いを細かく知っている必要はないと思うが,説明しておくと,\\
REF行列からpivotを1にし,pivotの上下を0にするとRREF行列となる.
長くなってしまったが\textbf{後退代入}のフェーズを説明する.
初心者向けにまず,前進消去後の行列を連立方程式に戻す(慣れればしなくていい).
\[
\begin{bmatrix}
    2& 4 & -2 & 2 \\
    0 & 1 & 1 & 4 \\
    0 & 0 & 4 & 8
\end{bmatrix}
\rightarrow
\begin{array}{ccccccc}
    2x & + &4y & -&2z &= &2 \\
       &  &1y & +&1z &= &4 \\
       &  &   &  &4z &= &8 \\
\end{array}
\]
見やすいように係数が0になる変数は消している.簡約化と異なり,第1行は変数が3つ生きています.しかし,「後退代入」をすると解ける.
前進消去は上から階段になるように,変数を消していく.恐らく,「前進消去」はここからきているのではないかと思っている.
同様に考えると「後退代入」はREF行列になった行列で下の行列から「代入」していく.つまり第3行から計算する.
第3行では変数が1つなので,$z=2$とすぐ求まる.次に第2行の計算をする.変数が2つだが,第3行で求めたzの値をzに「代入」する.
$y=2$と求まる.最後に第1行のx,yに「代入」すると,$x=-1$と求まった.\\
※この一連の処理を「ガウスの消去法(Gaussian Elimination)もしくは「掃き出し法(Row Reduction)」という.\\
※「前進消去(Forward Elimination)」「後退代入(Backward Substitution)」
\newpage
\paragraph{rankを求める}
rankはpivotの個数なのでREF行列の状態でも求められる. 以上\\
つまり,「簡約化」することが条件にされてなければ無理に簡約化する必要はないのである.\\
わからない？例を挙げよう.
\[
\begin{bmatrix}
    1& 2 & -1 & 1 \\
    4 & 9 & -3 & 8 \\
    -2 & -3 & 7 & 10
\end{bmatrix}
\]
こちらは先ほど,簡約化の手順の例でしようした行列である.簡約の手法と同様第1列から順番に処理していく.
\[
\begin{bmatrix}
    1& 2 & -1 & 1 \\
    0 & 1 & 1 & 4 \\
    0 & 1 & 5 & 12
\end{bmatrix}
\]
処理を終えた.簡約化では本来,2行2列の1の上下の要素を0にする.しかし階段化でに則り,第3行のみを0にする.
\[
\begin{bmatrix}
    1& 2 & -1 & 1 \\
    0 & 1 & 1 & 4 \\
    0 & 0 & 4 & 8
\end{bmatrix}
\]
ここで実は$rank$は3と求まってしまう.

第2行のpivotを用いて,第1行の2を0にした場合,変わる値は第1行の2,3,4列の要素である.$rank$はpivotの数というのであれば,果たしてここの数字が変わったことでpivotに影響があるのだろうか？
例えば,第1行のpivotは変わらない.第2行のpivotは処理をする側であり,される側ではないので変わらない.第3行のpivotが関わりそうな第1行3列の値が簡約化しないことで変わらなくても,同様に処理する側なので,影響を受けない.
したがって,\textbf{$\textbf{rank}$を求めるだけであれば簡約化までの処理は不要なのである!}
\subsubsection{列空間}\label{subsub:colspace}
カリキュラムが変わらなければ,線形代数学は「空間」というものに入っていく.これは空間のウォーミングアップになるといいと願う.\\
私が学習した内容で$A\symbfit{x}=\symbfit{b}$の考え方を変えてくれた素晴らしい内容である.\\
線形結合から行列にシフトしたタイミングを思い出して欲しい.$A\symbfit{x}$はAを構成する列ベクトルをそれぞれx倍y倍...としたベクトルたちをイメージする.
そのベクトルの合成によって指すベクトルが$\symbfit{b}$と一致することを式で表すと,$A\symbfit{x}=\symbfit{b}$となる.\\
では$A\symbfit{x}$が表現できるベクトルに限界はあるのだろうか？わかりやすい例としては,x倍y倍が自然数に限定した場合,明らかに,表現できない領域が現れる.\\
このように,Aの列ベクトルが表現できる領域を「列空間」という.\\
\begin{tcolorbox}[colframe=black!80,colback=white,title={$A\symbfit{x}=\symbfit{b}$における解の存在条件}]
$A\symbfit{x}=\symbfit{b}$が解を持つのは,$\symbfit{b}$が列空間に含まれるときであり,その時に限る.
\end{tcolorbox}
行列$A$の列空間はC(A)と書く.Cは列の英語の「Column」かと.
では列空間において基本的にx,yは実数全てを取るので,表現できない領域なんてあるのだろうか?と感じる.
列空間は名前からすると,3次元以上の一般的に「空間」と言われるものをイメージすると思うが,Aが平面を作る場合でもそれは列空間となる.
\newpage
具体的には
\[
C(A)=
x
\begin{bmatrix}1 \\2 \\3\end{bmatrix}
+y
\begin{bmatrix}2 \\ 1 \\ 1\end{bmatrix}
\]
は3次元の空間に平面を作る.この平面も列空間である.実際にやってみて欲しいが,自分がいる部屋で,紙を手で宙に浮かせて欲しい.$A\symbfit{x}=\symbfit{b}$を満たす場合,
この無限に広がるが限られた薄い平面上に基点からの$\symbfit{b}$が存在するということになる.なかなかシビアではないだろうか？※Aが直線を作る場合でも列空間という.\\
このように$A\symbfit{x}=\symbfit{b}$を満たすのは意外と大変である.そう考えると試験で出題される列空間はなかなか"作られている".現実では収まらないのことが多いらしく,
$\symbfit{b}$が列空間上に存在しないので,列空間上に存在する$\symbfit{b}$に一番近いベクトルを\textbf{近似解}として用いることがある.一番有名なのは\textbf{回帰直線}であるが,この物語の始まりはSection.\ref{sec:startGram}であるのでお楽しみに.
列空間の紹介はこの辺にしておく.\\

本来したかったrankの話をしよう.列空間はベクトルによって作られる空間であるが,その構成するベクトルはなるべく少ない方が情報量が少なくて嬉しい.
例えばこのような列空間はどうだろう?
\[
C(A)=
x
\begin{bmatrix}1 \\2 \\3\end{bmatrix}
+y
\begin{bmatrix}2 \\ 4 \\ 6\end{bmatrix}
+z
\begin{bmatrix}2 \\ 1 \\ 1\end{bmatrix}
\]

この列空間は3つの方向にベクトルが広がって3次元空間を作るように見えないだろうか？答えは間違いであり,正解は平面を作る.\\
というのもy倍されている列ベクトルはx倍されている列ベクトルの2倍であり,実は同じ方向を指している.つまり,3つの矢印は2つ重複しているので平面となる.
この考え方でいくと,列空間を表す列ベクトルは無限に増やせてしまい冗長である.言い換えれば,C(A)に必要なのは
$\begin{bmatrix}1 \\ 2\\ 3\end{bmatrix}$と$\begin{bmatrix}2 \\ 1\\ 1\end{bmatrix}$だけで十分であり,この2つでC(A)を十分表現できる.
ではrankと何の関わりがあるのか？実はrankは本来その行列が作る空間を表現する時に\textbf{最低何個の列ベクトルが必要か}を示してくれる.なので,この場合行列Aの$rank$は2となる.実際に調べてみても良い.「そうなんだ」「確かになる」と思っても,「何故」は残るだろう.
申し訳ないが今はこの程度の情報しか伝えることができない.しかし,これはかなり重要な情報であり,「独立と従属(Section.\ref{sec:depends})」という内容に入った時は必ず学べ,ここに戻ってくる.
最終的にはpivotとrankと「独立と従属」がSection.\ref{sub:space_base}で完全に繋がるだろう.小さな豆知識程度のものだとは思わないでほしい.
\newpage
\subsection[\texorpdfstring{\textcolor{blue}{(数値解析)数値線形代数学}}{(数値解析)数値線形代数学}]{(数値解析)数値線形代数学}
このセクションでは24年度「数値解析」に関する内容と先程まで説明していた内容に関することも含め「数値線形代数学」の話をする.数値解析の内容に興味を持った方はこの内容に対して面白く読むことができると考えている.
\subsubsection[\texorpdfstring{\textcolor{blue}{最適化された高速行列・ベクトル計算ライブラリ}}{最適化された高速行列・ベクトル計算ライブラリ}]{最適化された高速行列・ベクトル計算ライブラリ}
例えばガウスの消去法をプログラムで作成した場合,高速性はコード作成の工夫のみならず,演算するパソコンの処理速度自体が関わっている.例えば,2008年5月20日にIBMとロスアラモス研究所から,スパコン「Roadrunner」の新しい記録が発表されました.このスパコンは
1秒間に1000兆回($10^{15}$回)の浮動小数点演算を最初に達成しました.では最近はどうだろうか？2024年に米国ローレンス・リバモア国立研究所のスパコン「El Capitan」が1秒間に$10^{18}$回の計算を達成している.

$A\symbfit{x}=\symbfit{b}$の解を解くために,長々とガウスの消去法を組むのは現実的ではない.Python3のNumpyなどのモジュールで1行で解けてしまうようなものが用意されている.
別にこれは魔法を呼ぶコマンドではなく,最適化された\textbf{BLAS}(\textit{Basic Linear Algebra Subroutines})によって事前に組まれたプログラムが高速演算を呼び出し内部で行い,その結果を出力している.
BLASは低レベルな基本的な線形代数演算を担い,これ以降に紹介する固有値やQR分解などはBLASの上位レイヤーとなる\textbf{LAPACK}(\textit{Linear Algebra PACKage)})が担当している.この\underline{内部処理に使われる「部分ピポット選択」についてSection.\ref{subsub:partial_pivot}少し触れる}予定である.

高度に最適化されたBLASでは,行列の計算についての考え方がうまく反映されている.BLASにはレベル1-3がある.
\begin{enumerate}[label=]
    \item \textbf{レベル1} ベクトルの線形結合$a\symbfit{u}+\symbfit{v}$：計算量$O(n)$
    \item \textbf{レベル2} 行列--ベクトル積$A\symbfit{u}+\symbfit{v}$：計算量$O(n^2)$
    \item \textbf{レベル3} 行列--行列積$AB+C$：計算量$O(n^3)$
\end{enumerate}
\subsubsection[\texorpdfstring{\textcolor{blue}{命令数から見るGaussian Eliminationの利点}}{命令数から見るGaussian Eliminationの利点}]{命令数から見るGaussian Eliminationの利点}\label{subsub:Gaussian_order}
Section.\ref{sub:Gaussian}では前進消去と後退代入の一連の操作について紹介した.ここではこのガウスの消去法アルゴリズムの「\textbf{命令数}」について触れようと思う.「計算量」は$O(n)$など個々の命令の数は気にせず,入力サイズに応じてどれくらいの計算が増えるのかという\textbf{増加の傾向}を見るためにある.
対して,「命令数」は\textbf{実行する命令の総数}を意味する.例えば,for文がn回の内,加算が4回行われた場合,「計算量」は$O(n)$となり,「命令数」は$n+4$となります.

\subsubsection[\texorpdfstring{\textcolor{blue}{消去の破綻と部分ピボット選択}}{消去の破綻と部分ピボット選択}]{消去の破綻と部分ピボット選択}\label{subsub:partial_pivot}

\subsubsection[\texorpdfstring{\textcolor{blue}{丸め誤差と完全ピボット選択}}{丸め誤差と完全ピボット選択}]{丸め誤差と完全ピボット選択}\label{subsub:full_pivot}

\newpage
\section{(重要)行列操作の規則}\label{sec:matrixRule}
\subsection{和と差とスカラー倍の法則}
$A\symbfit{x}=\symbfit{b}$で行列同士の掛け算を少しやっているが,改めてまとめる.これから行列同士の計算が本格的に必要になってくるからだ.
\begin{tcolorbox}[colframe=black!80,colback=white,title={行列和と行列差の計算可能な条件,スカラー倍の計算}]
和と差について,\textbf{行列の\hyperlink{size}{サイズ}が同じであれば}計算することができる.
\[
\begin{bmatrix}
    1 & 2 \\ 3 & 4 \\ 0 & 0
\end{bmatrix}
+
\begin{bmatrix}
    2 & 2 \\ 4 & 4 \\ 9 & 9
\end{bmatrix}
=
\begin{bmatrix}
    3 & 4 \\ 7 & 8 \\ 9 & 9
\end{bmatrix}
\]
また,行列に対して何倍かする時,その係数は\textbf{すべての要素に施される}.
\[
2\begin{bmatrix}
    1 & 2 \\ -3 & 4 \\ 0 & 5
\end{bmatrix}
=
\begin{bmatrix}
    2 & 4 \\ -6 & 8 \\ 0 & 10
\end{bmatrix}
\]
\end{tcolorbox}
Aを-Aに足すと,すべての要素は0である「\textbf{零行列}」となる.\\
簡単だと感じただろうか？しかし,本格的な問題は「行列積」.つまり行列同士の掛け算である.\\
これからそれについて触れる前に要素の示し方について触れておこう.i行j列にある要素は$a_{ij}$また$A(i,j)$と書くのが一般的である.
例えば,第1行のn個の要素は$a_{11},a_{12}, \cdots ,a_{1n}$となる.行列の左下の要素は$a_{m1}$となり,右下は$a_{mn}$となる.
\subsection{積の法則}
\subsubsection{要素同士の掛け合わせ手順}
では行列積の話をしよう.行列積には和と差同様,行列の形によっては計算できないときがある.
\begin{tcolorbox}[colframe=black!80,colback=white,title=行列積の計算可能なサイズの組み合わせ]
行列$A$,Bが計算できるためには,「Aがn列からなるとき,Bはn行からなる必要がある」\\
文字で見ると難しいが視覚的にわかりやすい書き方がある.
\\$A:m\times n \; \; B:n\times k$\\ 
の時,隣り合う数字が同じなら計算可能\\
この場合,nのことである.
また,行列積の結果は,隣り合う数字が消されて結合されたように,$m\times k$となる.\\
\end{tcolorbox}
例えば,
\[
\begin{bmatrix}
    1 & 2 \\ 3 & 4 \\ 0 & 0
\end{bmatrix}
\begin{bmatrix}
    2 & 2 \\ 4 & 4 \\ 9 & 9
\end{bmatrix}=?\qquad(\text{※行列同士の積は}\times\text{や}\cdot\text{を書かない})
\]
は計算ができない.何故なら,$3\times2$と$3\times2$で隣り合う数字が異なるからである.
しかし,
\[
\begin{bmatrix}
    1 & 4 \\ 1 & 5 \\ 4 & 2
\end{bmatrix}
\begin{bmatrix}
    3 & 2 \\ 1 & 0
\end{bmatrix}
\]は計算が可能である.サイズを見ると,隣り合う数字が2であるからである.

計算方法であるが,まずは授業で紹介される書き方をまず真似てみるといい.以降の文章は\underline{それを習った前提で}進める.

私が勝手に思っているかもしれないが,計算結果を書いていく時,掛け合わせのミスは無くなるだろうが,純粋な和の結果を描くときに計算ミスをする可能性があるなと感じている.
和の計算に限っては暗算になるからである.問題点は和が正しいか検算する時に途中式がないので,再度計算し直す必要があるということである.小さい結果になるならそんな話はしないが,分数同士や2桁になってくるとまずい.時間がかかる.
そんなこと気にする人の方が少ないかもしれないし,
気にしない人は無視していいのだが,計算ミスが多い私としては工夫したいところではある.\medskip\\
\fbox{ポイント}
\begin{framed}
列\times 行という方法でも行列積は求まる.\\
\[
\begin{bmatrix}
    1 & 4 \\ 1 & 5 
\end{bmatrix}
\begin{bmatrix}
    3 & 2 \\ 1 & 0
\end{bmatrix}
=
\begin{bmatrix}
    1 \\ 1
\end{bmatrix}
\begin{bmatrix}
    3 & 2
\end{bmatrix}
+
\begin{bmatrix}
    4 \\ 5
\end{bmatrix}
\begin{bmatrix}
    1 & 0
\end{bmatrix}
=
\begin{bmatrix}
    3 & 2 \\ 3 & 2
\end{bmatrix}
+
\begin{bmatrix}
    4 & 0 \\ 5 & 0
\end{bmatrix}
=
\begin{bmatrix}
    7 & 2 \\ 8 & 2
\end{bmatrix}
\]
授業でやった計算の組み合わせとは本質的には変わっていない.
\end{framed}
計算がシンプルに見え,かつ工程が多く見れるので検算しやすいのではと感じている.これは好みの問題だが,私はこののやり方がやりやすい.
さらにこれは\textbf{サイズの組み合わせ}に従っている.

しかし,\textbf{必ず注意して欲しいのは,行\times 列はではない.}ただ,実際にこれは計算すると誤っていることに気づく.
\underline{行列のサイズによらず,行\times 列をすると必ず,結果は1\times1の行列となる.}理由はシンプルでサイズが大きくなっても長い列ベクトル,短い列ベクトルとなるだけで,列ベクトル同士の掛け算において行と列どちらかは必ず1となるからである.
\subsubsection{行列同士の掛け合わせ手順}
上でその行列同士の積が取れるかどうかはサイズによると説明した.そのため,今までやってきたような$3\times4=12$のような掛け算で行えた手法が使えなくなる.このセクションではそれについて説明する.短い内容だが,\textbf{行列同士の計算の根底となるため非常に重要である}.
行列$A,B,C$において,
\[
    AB=CD
\]
が成りたっているとする.ただし,$A$のサイズは$n\times k$とする.\\
このとき,$m\times n$の行列$E$を掛けるとき,サイズ的に$EA$と積はとれる.
\[
    EAB=ECD
\]
と計算可能である.では,$3\times4=12$のように順番を並び替えて,
\[
    EAB=CED,\quad EAB=CDE
\]
などとしていいのだろうか？答えはバツである.理由は単純で,サイズの組み合わせで計算できない\underline{可能性がある}からである.例えば,行列$D$のサイズが$j \times p$なら$p= m$とは限らない.
\[
\text{行列を両辺で掛けるときは,\textbf{左から},\textbf{右から}である.}
\]

\[
AB = C \quad\Longrightarrow\quad
\begin{cases}
EAB = EC\\
ABE = CE
\end{cases}\\
\text{※いずれもサイズ的に計算可能と仮定する}
\]

では,\underline{可能性がある}と書いたが,サイズの組み合わせによって計算できたら順番を並び替えてもいいのか？と思うかもしれない.
実はそうではない.入れ替えしても問題ない行列は$AB=BA$が成り立つ行列に限り可能であり,これを可換という.可換が成り立つ代表例はあとに紹介する逆行列かもしれない.しかし,ここまで学ぶ必要は無いと思うし,可換が成り立つかもわからない行列に対して無理に入れ替えして間違えるのは危険である.
\section{逆行列}
\subsection{逆行列とは}
逆行列は習っただろうか？計算手法は次のSection.\ref{sub:Gauss_jordan}で行っているが,ここでは逆行列とはそもそもなんなのか見てみよう.
逆行列は求めることがテストに出題されるが,求めて何になるのだろうか？と感じる人も一定数いるかもしれない.

行列$A$の逆行列は$A^{-1}$と書き,「Aの逆行列」と読んでもいいが,私は「エーインバース」と読んでいる.
\begin{tcolorbox}[colframe=black!80,colback=white,title=逆行列の重要性質]
行列$A$に対する逆行列$A^{-1}$とは,
\[AA^{-1}=A^{-1}A = I\]
が成り立つ行列である.ただし,Iは単位行列である.※単位行列は対角成分が1でそれ以外の要素が全て0な正方行列(n\times n行列)
\end{tcolorbox}
積をとったら単位行列になるという性質は逆行列の形や逆行列を何かの行列にかけるとき面白い性質を見せる.Section.\ref{sub:I_inv}(発展)で紹介する.\\
また,求めた逆行列が計算ミスを起こしていないかの確認の手段に使える.正しく逆行列が求められていれば単位行列になるはずである.
私は名称好きなのかもしれない.逆行列を求める方法にも名前があるので紹介する.以下のタイトルが逆行列を求める方法である.
\subsection{(重要)ガウス-ジョルダンの消去法(Gauss-Jordan Method)}\label{sub:Gauss_jordan}
具体的な方法の前になぜ求まるかだけ説明させてほしい.\\
Aは3次正方行列として解説する.求める$A^{-1}$の列ベクトルを$x_1,x_2,x_3$とし,単位行列の列ベクトルを$e_1,e_2,e_3$とする.
\[
AA^{-1}=A\begin{bmatrix}x_1 & x_2 & x_3\end{bmatrix} = \begin{bmatrix}e_1 & e_2 & e_3\end{bmatrix} = I
\]
となる.
これより拡大係数行列$\begin{bmatrix} A & I\end{bmatrix}$を用意する,この行列に$A^{-1}$をかけたとしたら,
$\begin{bmatrix} I & A^{-1}\end{bmatrix}$になることは何となく上の性質からわかるだろう.\\
つまり3次正方行列の逆行列を求めたい場合は,AとIを横並びにした3\times 6の拡大係数行列を用意して,左半分が単位行列になるように,
簡約化と同じ手法を取れば右半分はAの逆行列となる.\\
試しに,$A=\begin{bmatrix}
    2 & -1 & 0 \\
    -1 & 2 & -1 \\
    0 & -1 & 2
\end{bmatrix}$の逆行列を求めてみる.\\
※簡約化のやり方は説明しているので,ちょっと違ったプロセスを辿って書いてみる.
\[
\begin{bmatrix}A & e_1 & e_2 & e_3 \end{bmatrix}=
\begin{bmatrix}
2 & -1 & 0 & 1 & 0 & 0 \\
-1 & 2 & -1 & 0 & 1 & 0 \\
0 & -1 & 2  & 0 & 0 & 1
\end{bmatrix}
\rightarrow
\begin{bmatrix}
2 & -1 & 0 & 1 & 0 & 0 \\
0 & \frac{3}{2} & -1 & \frac{1}{2} & 1 & 0 \\
0 & -1 & 2  & 0 &0 & 1
\end{bmatrix}
\rightarrow
\begin{bmatrix}
2 & -1 & 0 & 1 & 0 & 0 \\
0 & \frac{3}{2} & -1 & \frac{1}{2} & 1 & 0 \\
0 & 0 & \frac{4}{3} & \frac{1}{3} & \frac{2}{3} & 1
\end{bmatrix}
\]
\[
\rightarrow
\begin{bmatrix}
2 & -1 & 0 & 1 & 0 & 0 \\
0 & \frac{3}{2} & 0 & \frac{3}{4} & \frac{3}{2} & \frac{3}{4} \\
0 & 0 & \frac{4}{3} & \frac{1}{3} & \frac{2}{3} & 1
\end{bmatrix}
\rightarrow
\begin{bmatrix}
2 & 0 & 0 & \frac{3}{2} & 1 & \frac{1}{2} \\
0 & \frac{3}{2} & 0 & \frac{3}{4} & \frac{3}{2} & \frac{3}{4} \\
0 & 0 & \frac{4}{3} & \frac{1}{3} & \frac{2}{3} & 1
\end{bmatrix}
\rightarrow
\begin{bmatrix}
1 & 0 & 0 & \frac{3}{4} & \frac{1}{2} & \frac{1}{4} \\
0 & 1 & 0 & \frac{1}{2} & 1 & \frac{1}{2} \\
0 & 0 & 1 & \frac{1}{4} & \frac{1}{2} & \frac{3}{4}
\end{bmatrix}=
\begin{bmatrix}
    I & A^{-1}
\end{bmatrix}
\]
このやり方は手計算では辛いだろうから,簡約化の方法に従った方がいい.手を動かして,結果が同じなることを確認せよ.ただ,この過程を辿ったの理由は「\textbf{行列式}」を習う時に使いたいからだ.
\subsection{(発展)単位行列から見る逆行列の性質}\label{sub:I_inv}
単位行列の性質として以下のものがある.
\[
I^nA=A
\]
つまり,単位行列をいくらかけても行列に変化が起きないのである.
ここから面白いものが見える.それは\textbf{行列$A$によっては逆行列$A^{-1}$はガウス-ジョルダンの消去法で計算しなくても求められる}ということである.
例えば,こんな行列
\[
 A=4I = \begin{bmatrix}
    4 & 0 & 0 \\
    0 & 4 & 0 \\
    0 & 0 & 4
\end{bmatrix}
\]
この行列をとある行列Bにかけると,Bのすべての要素は4倍になる.
\[
A\begin{bmatrix}
    2 & 3 \\ 5 & 1
\end{bmatrix}
=4I\begin{bmatrix}
    2 & 3 \\ 5 & 1
\end{bmatrix}
=4\begin{bmatrix}
    2 & 3 \\ 5 & 1
\end{bmatrix}
=
\begin{bmatrix}
    8 & 12 \\ 20 & 4
\end{bmatrix}\qquad
\left(
    B=\begin{bmatrix}
    2 & 3 \\ 5 & 1
\end{bmatrix}
\right)
\]
ではこのAの逆行列を考えてみよう.
逆行列の性質的に,
\[
A^{-1}AB = IB = B
\]
すなはちAによってBが4倍されたとしても,$A^{-1}$によって元の行列$B$に戻るのが逆行列となりそうである.
元に戻すにはBのすべての要素を$\frac{1}{4}$すればいいので,
\[
A^{-1} = \begin{bmatrix}
    \frac{1}{4} & 0 & 0 \\
    0 & \frac{1}{4} & 0 \\
    0 & 0 & \frac{1}{4}
\end{bmatrix}
\]
とガウス-ジョルダンの消去法をしなくても求められることがある.
他にも転置行列や置換行列,基本行列などでも同様の性質を活用して計算せずとも求めることができる場合がある.\\
例えば,基本行列の逆行列を理解すると,後に習う「LU分解」を手計算で求めることができる.

\section{(重要)「拡大係数行列の簡約行列」から連立一次方程式を解く}
前チャプターの「線形方程式の求解」でのつづきてである.つまり,フライングしている.その問題で求まる解はただ1つに定まっていた.しかし,解が1つに定まらない場合はすでに習っている.
媒介変数とか言っただろうか？ここではパラメータという名前で説明していく.このチャプターでは線形方程式の解き方がより一般的になる.これは1Qの出題範囲でもあり,2Qの出題範囲でもある.
また,これはrankとpivotが深く関わる.忘れた人は復習してから読み進めてほしい.
基本的な計算方法はすでに習得している.拡大係数行列に対して,簡約化を行う.これは解がどのようになろうとも同じである.
そこからどのような解になるかはpivotの位置と個数がそれを決める.つまり解がどのようになるかは簡約化してから知れる物であり,パラメータになろうとも
工程はほとんど変わらないため覚えることもそれほど多くない.

では以下の問題を解いてみよう.\\
$\begin{cases}
    2x-y+z=4 \\
    x+2y-z=3 \\
    2x-3y+2z=3
\end{cases}$
拡大係数行列を用いて簡約化まで行う.\\
$[A | \symbfit{b}]=
\begin{bmatrix}
    2 & -1 & 1 & 4 \\
    1 & 2 & -1 & 3 \\
    2 & -3 & 2 & 3
\end{bmatrix}
\rightarrow
\begin{bmatrix}
    1 & 2 & -1 & 3 \\
    0 & -5 & 3 & -2 \\
    0 & 6 & -3 & -3
\end{bmatrix}
\rightarrow
\begin{bmatrix}
    1 & 0 & 0 & 2 \\
    0 & 1 & 0 & 1 \\
    0 & 0 & 1 & 1 
\end{bmatrix}$\\
ここで再び,連立一次方程式に戻すと,
$\begin{cases}
    1x+0y+0z=2 \\
    0x+1y+0z=1 \\
    0x+0y+1z=1
\end{cases}$となり,これはすぐ解がわかる.\\
またもう1つわかることがある.それは解が1つに定まっているということである.\bigskip\\
ではこの問題はどうだろうか？\\
$\begin{cases}
    x+2y+2z=4 \\
    x+3y+5z=5 \\
    2x+y-5z=5
\end{cases}$\\
同様に行うと,\\
$[A | \symbfit{b}]=
\begin{bmatrix}
    1 & 2 & 2 & 4 \\
    1 & 3 & 5 & 5 \\
    2 & 1 & -5 & 5
\end{bmatrix}
\rightarrow
\begin{bmatrix}
    1 & 2 & 2 & 4 \\
    0 & 1 & 3 & 1 \\
    0 & -3 & -9 & -3
\end{bmatrix}
\rightarrow
\begin{bmatrix}
    1 & 0 & -4 & 2 \\
    0 & 1 & 3 & 1 \\
    0 & 0 & 0 & 0 \\
\end{bmatrix}$\\
連立一次方程式に戻すと,
$\begin{cases}
    1x+0y-4z=2 \\
    0x+1y+3z=1 \\
    0x+0y+0z=0
\end{cases} \rightarrow
\begin{cases}
    1x-4z=2 \\
    1y+3z=1 \\
\end{cases}
$となる.\\
結論としてはこれはパラメータ付きの解となる.(試してみるといい)\medskip\\
この違いはなんだろうか？\\
答えは求める解の個数に対して,\underline{方程式の個数が足りない}ということである.\\
1つの解が求まる最低条件は「方\textbf{程式の数=未知数の数}」ですので,これはそれを満たさない.\\
この場合,x,y,zのどれかに対してパラメータにすることで解を表現できることになるが,どれをパラメータにすれば良いのだろうか？
x=1にしてyを考える？実は簡約後の拡大係数行列のpivotに注目するとこの問題は比較的容易に解決することができる.
\vspace{1cm}


2つの話をしよう.
\begin{enumerate}
    \item 係数行列$A$と解からなる列ベクトル$\symbfit{x}$を見ると,Aの各列の要素に変数が掛けられていく結果が方程式となる.解がx,y,zとするなら,Aの第1列にはx,第2列にはyという感じに.
    \item A全体で見ると,pivotがある列とない列が存在する.行列の簡約化時にはpivotの値が前後の行にある要素を変えることで,簡約化が進んでいる.
        pivotが定まってからpivotは常に変わらず,pivotではない値が変化する.簡約化したAは誰が行っても同じ結果になるが過程は異なる.簡約化は実行前から結果が確定していることが言える.
\end{enumerate}
つまり,\\
「\textbf{pivotは簡約化前からすでに値が確定している定数で,それ以外の要素は変数である.}」と言えないだろうか？\\
1つ目の話と私の言いたいことが見えてくる.では第1列にpivot(行は自由)が存在する場合,それに掛けられるxは同様に定数になると考えるのは自然ではないだろうか？\\
列においてpivotは1つだが,連立されているどの方程式は同じ解を示す.どれか1つの方程式でpivotと解が掛けられていれば,別の方程式では掛けられていないくても定数であることは変わらない.
1行1列にpivotがある場合,xとpivotが掛けられることになる.xの解は定数となる.第2,3行の第1列にはpivotはないが,3つの方程式は同じ解を指すので,xが定数となる.\\
これを踏まえて具体例を見よう.\\

$\begin{cases}
    2x-y+z=4 \\
    x+2y-z=3 \\
    2x-3y+2z=3
\end{cases}$
は簡約化により,
$\begin{cases}
    1x+0y+0z=2 \\
    0x+1y+0z=1 \\
    0x+0y+1z=1
\end{cases}$となった.\\
pivotの位置に注目すると,すべての列に存在する.つまり,すべての解x,y,zに対応するpivotが掛けられることを意味する.
確かに解が1つ.つまりすべて定数の解を導き出している.\\

パラメータが出現する問題を見てみよう.\\
$\begin{cases}
    x+2y+2z=4 \\
    x+3y+5z=5 \\
    2x+y-5z=5
\end{cases}$は拡大係数行列の状態では
$
\begin{bmatrix}
    1 & 0 & -4 & 2 \\
    0 & 1 & 3 & 1 \\
    0 & 0 & 0 & 0 \\
\end{bmatrix}$でzが掛けられる列にpivotがない.\\
つまり,zが変数(パラメータ)になることがわかる.\\
この問題の最終ゴールを示すなら,zがパラメータよりz=tとする.
それより,x,yをtを用いて表現すると,$x=t+2,y=-3t+1$となる.解を線形結合の形で書くと,\\
\[
\begin{bmatrix}
    x \\ y \\ z
\end{bmatrix}=t
\begin{bmatrix}
    4 \\ -3 \\ 1
\end{bmatrix}+
\begin{bmatrix}
    2 \\ 1 \\ 0
\end{bmatrix}
\]\\
これで求められた.\\

\begin{tcolorbox}[colframe=black!80,colback=white,title= パラメータの存在条件]
\begin{center}
   \textbf{行列Aの列数=pivotの個数}
\end{center}
    \hspace{1cm}$\Longrightarrow\:$\textbf{pivotがない列がパラメータとなる.}\\
    \hspace{1cm}$\Longrightarrow\:$\textbf{pivotがない列数とパラメータの数が一致する.}\\
    \hspace{1cm}$\Longrightarrow\:$\textbf{全ての列にpivotがある場合,解が一意に定まる.(パラメータが存在しない)}
\end{tcolorbox}
特に最後の解が一意に定まる行列を「\textbf{非退化}な行列」という.これはSection.\ref{sub:nondegenerate}で扱うため,読んだ後に振り返るとより理解が深まるかもしれない.
\section{行列式について}\label{sec:det}
(注意！！)行列式の内容に入ってから読むことを前提とする.\\
行列式の内容は今までの学習内容と異なり,計算メインとなっている.そのため特段説明することがないというのが本音ではある.\\
授業では「サラスの公式」や「余因子展開」がなぜ成り立つかについてはこのテキストでも紹介しない.発展で扱うべきであっても,
それを知ったことによる応用方法を知らないからである.例えば,今私が「余弦定理はベクトルから証明できるから示す.」といって書いたとしても
「へー」で終わってしまう.\\
ただ,同じタイミングで逆行列と可逆・非可逆についての内容に触れているだろうから,それについてまとめる.
また,(発展)では「なぜ」について触れる.これは知っておくとテストにおいての点数の差はないもののこれから先の
線形代数学や関連内容の見通しが良くなるだろうからおすすめする.
\subsection{逆行列と可逆性}
まず前提として「正則な行列」とはなんなのだろうか？\\
線形代数学では「正則」と「可逆」は同じ意味である.理解しやすい「可逆」を使うが,テストに向けて最後は「正則」に自己変換して覚えてほしい.\\
再びAx=bで説明する.Aの逆行列を用いると,$\symbfit{x}=A^{-1}\symbfit{b}$と解を表現することができる.
いままでの数学ではパラメーターが出現することを見かけたことがあるかと思う.同様に行列$A$によっては解が
1つに定まるとは限らない.解がパラメーターとなることはある.\\
例えば$A\symbfit{x}=\symbfit{b}$において,この式を満たす$\symbfit{x}$が2つあったとする.では$\symbfit{x}=A^{-1}\symbfit{b}$が表す解は同じく2つである.2つの式はイコールなので同値変形となる.
同値変形は式変形をした時の前後で結果が変わってはいけない.式変更後と前で常に「どちらかの解」となってしまっているので同値変形とはならない.
つまり,$\symbfit{x}=A^{-1}\symbfit{b}$が成り立つ時,解が1つでなければならない.
これを転じて言えることは,
\begin{center}
    「行列$A$が逆行列を持つ時,$A\symbfit{x}=\symbfit{b}$を満たす解はただ1つでなけらばならない.」    
\end{center}
このような行列を「可逆な行列」という.逆を「非可逆な行列」と言う.\\

先程の説明によって$\symbfit{b}$の値によって$A\symbfit{x}=\symbfit{b}$の解が1つかそれ以外か決まるのではないかと思ったかもしれない.\\
実は可逆かどうかは行列$A$のみで決まる.(発展)で説明する.

\subsubsection{行列式と逆行列の可否}
では行列が可逆であるかを判断する方法がもちろん存在する.\\
(重要)
\begin{center}
    $\det(A) \neq 0$である時,行列$A$は可逆である.
    それはつまり,逆行列$A^{-1}$をもつ.
\end{center}
なぜそのようになるのかについては(発展)で説明するが,この行列式が逆行列をもつか？や可逆か？
と問われたら行列式を計算する.\\
また,次のことが言える.
「逆行列を持つ時,必ず解を持つ.」なぜなら確実に$\symbfit{x}=A^{-1}\symbfit{b}$とできるので,$\symbfit{x}$の存在を約束することができるからだ.
逆行列をもたない行列の場合は「解なし or 複数解」となる.


\subsection{(発展)前進消去から見る可逆と非可逆と逆行列と}
ここでは先ほどの説明の「なぜ」に注目していこう.この説明に向けて伏線があったので予言通り回収しておく.
\[
\begin{bmatrix}A & e_1 & e_2 & e_3 \end{bmatrix}=
\begin{bmatrix}
2 & -1 & 0 & 1 & 0 & 0 \\
-1 & 2 & -1 & 0 & 1 & 0 \\
0 & -1 & 2  & 0 & 0 & 1
\end{bmatrix}
\rightarrow
\begin{bmatrix}
2 & -1 & 0 & 1 & 0 & 0 \\
0 & \frac{3}{2} & -1 & \frac{1}{2} & 1 & 0 \\
0 & -1 & 2  & 0 &0 & 1
\end{bmatrix}
\rightarrow
\begin{bmatrix}
2 & -1 & 0 & 1 & 0 & 0 \\
0 & \frac{3}{2} & -1 & \frac{1}{2} & 1 & 0 \\
0 & 0 & \frac{4}{3} & \frac{1}{3} & \frac{2}{3} & 1
\end{bmatrix}
\]
\[
\rightarrow
\begin{bmatrix}
2 & -1 & 0 & 1 & 0 & 0 \\
0 & \frac{3}{2} & 0 & \frac{3}{4} & \frac{3}{2} & \frac{3}{4} \\
0 & 0 & \frac{4}{3} & \frac{1}{3} & \frac{2}{3} & 1
\end{bmatrix}
\rightarrow
\begin{bmatrix}
2 & 0 & 0 & \frac{3}{2} & 1 & \frac{1}{2} \\
0 & \frac{3}{2} & 0 & \frac{3}{4} & \frac{3}{2} & \frac{3}{4} \\
0 & 0 & \frac{4}{3} & \frac{1}{3} & \frac{2}{3} & 1
\end{bmatrix}
\rightarrow
\begin{bmatrix}
1 & 0 & 0 & \frac{3}{4} & \frac{1}{2} & \frac{1}{4} \\
0 & 1 & 0 & \frac{1}{2} & 1 & \frac{1}{2} \\
0 & 0 & 1 & \frac{1}{4} & \frac{1}{2} & \frac{3}{4}
\end{bmatrix}=
\begin{bmatrix}
    I & A^{-1}
\end{bmatrix}
\]
これは逆行列の求める方法で紹介された内容をコピペしてきたものである.前ページの説明より,逆行列と行列式には関係があることがわかるので,この変形にも行列式の一端が見れる.\\
注目するのは変形過程の最後から2番目の行列である.対角に$2,\frac{3}{2},\frac{4}{3}$と並んでいる.これはpivotですよね.\\
この積を「pivot積」というが,
\begin{center}
  \bf  「pivot積」=$\: \det(A)$
\end{center}
である.つまり,$\det(A)= 2 \cdot \frac{3}{2} \cdot \frac{4}{3}=4$と求めることが可能である.\\
ここから急速に点と点が結ばれていく.\\
ここで注意してほしいのは必ず成り立つわけではない.行列式の性質を思い出してほしい.行列式では行の定数倍は約数として,外にだすことが可能であった.また,行(列)を交換すると
マイナスがつく性質があった.そのため,この対角線上にpivot
が現れるこの変形はそれに注意を払った変形であることに気をつけてほしい.\\
では行列式ではなく$A\symbfit{x}=\symbfit{b}$を解いていると仮定してよう.\\
$\symbfit{b}$の値によらず,下記のようになっていくと思う.\\
$A\symbfit{x}=\symbfit{b} \rightarrow$
$\begin{bmatrix}
    2 & 0 & 0 \\
0 & \frac{3}{2} & 0 \\
0 & 0 & \frac{4}{3}
\end{bmatrix}
\begin{bmatrix}
    x \\ y \\ z
\end{bmatrix}=
\symbfit{b'}$\\
この時,$x=\frac{1}{2}b'_1,...$と解が求まっていく.ではこの対角の数字のうち1つが0の場合を考えてほしい.そうなるとrank(A)=2となり,解がパラメータ付きの「無数解」になるか$0x=b'_1(\neq 0)$で「解なし」になる.
この時の「pivot積」は0となる.それはすなわち$\det(A)=0$である.これが$\det(A) \neq 0$の時,解が1つであり,=0の時「解なし」「無数解」となる理由である.\\

次に逆行列である.\\\\
$\begin{bmatrix}
    2 & 0 & 0 & & & \\
    0 & \frac{3}{2} & 0 & & J & \\
    0 & 0 & \frac{4}{3} & & &  
\end{bmatrix}
\rightarrow
\begin{bmatrix}
    1 & 0 & 0 & & \frac{1}{2}(Jの第1行) & \\
    0 & 1 & 0 & & \frac{2}{3}(Jの第2行) & \\
    0 & 0 & 1 & &  \frac{3}{4}(Jの第3行)&  
\end{bmatrix}
\rightarrow
\begin{bmatrix}
    1 & 0 & 0 & & & \\
    0 & 1 & 0 & & \frac{1}{\det(A)}J & \\
    0 & 0 & 1 & & &  
\end{bmatrix}
\quad \therefore
\frac{1}{\det(A)}J = A^{-1}$\\\\
となる.\underline{(分母)$\neq 0$であるから$\det(A) \neq 0$でないと逆行列が存在しない}ことがわかる.
\section{外積と内積}\label{sec:inner_outer}
\subsection{転置を用いた内積の表現}
\[
\symbfit{a}=\begin{bmatrix}
    1 \\ 2 \\ 3
\end{bmatrix},\qquad \symbfit{b}=\begin{bmatrix}
    -3 \\ -2 \\ -1
\end{bmatrix}
\]
\begin{align}
  \symbfit{a}\cdot \symbfit{b}&=1\cdot(-3)+2\cdot(-2)+3\cdot(-1)=0\label{eq:inner1}\\
  \symbfit{a}^T\symbfit{b}&=\begin{bmatrix}1 & 2 & 3\end{bmatrix}
\begin{bmatrix}-3 \\ -2 \\ -1\end{bmatrix}=1\cdot(-3)+2\cdot(-2)+3\cdot(-1)=0\label{eq:inner2}
\end{align}

高校数学において,内積は式(\ref{eq:inner1})のように$\vec{a}\cdot\vec{b}$と表現していた.しかし,私たちは\textbf{転置}を学んだので,\underline{転置を用いた内積表現が可能となる.}
私のこだわりだが,行列の演算をしているなら,内積も行列の演算で表現したい.
行列の計算では,AB なら A の第 1 行と B の第 1 列というように十字に計算する.行列の計算法則で内積の計算をするにはどうしたらいいだろうか？それは転置である.列ベクトル AB の内
積を計算したい場合,A を転置させることで,縦に長い行列から横に長い行列に変わる.これが式(\ref{eq:inner2})である.このまま行列の積を取れ
ば,実質今までの内積同様計算方法に帰着することができる.
\begin{tcolorbox}[colframe=black!80,colback=white,title=転置表現での内積$_{\text{inner product}}$]
\[
\symbfit{a}\cdot\symbfit{b} = \symbfit{a}^T \symbfit{b}
\]
\end{tcolorbox}
転置表現のいいところは多数のデータで一気に内積を取れるところである.
今までの書き方では行列(列ベクトルではない)のような平面的な数字の羅列に対する内積表現ができなかった.
他にもメリットがある.\\式(\ref{eq:inner1})からわかるように,内積を表すドットは掛け算と似ている上,小さくて見えずらい.読者は積極的に転置表現を使って欲しい.

特に内積 0 は重要な意味を果たす.
これは工学や経済学など幅広い分野でつかわれている.Section.\ref{sec:ad_inner}で特集している.
\subsection{外積}
\subsubsection{外積とは}
外積は一部の高校生は計算を簡単にするためのテクニックとして学ぶも方もいるが,習わない人が多い.この外積は内積と似た形をいているのに習わないのはおかしい.さらに言えば,
この外積は非常に便利である.高校数学で使われたように,特に高校数学で輝く.
ここでは外積とは何なのか説明していく.\bigskip\\

2つの零でないベクトル$\symbfit{a}=\begin{bmatrix}a_1 & a_2 & a_3\end{bmatrix}^T,\symbfit{b}=\begin{bmatrix}b_1 & b_2 & b_3\end{bmatrix}^T$が与えられるとき,$\symbfit{a}$と$\symbfit{b}$の両方に直交するベクトル$\symbfit{c}=\begin{bmatrix}c_1 & c_2 & c_3\end{bmatrix}^T$を求めるとする.
\begin{align}
   \symbfit{a}^T\symbfit{c}=a_1c_1+a_2c_2+a_3c_3=0\\
   \symbfit{b}^T\symbfit{c}=b_1c_1+b_2c_2+b_3c_3=0
\end{align}
$c_3$を削除するために,式(1)の両辺を$b_3$倍,式(2)の両辺を$a_3$倍して,引き算すると,
\[
    (a_1b_3-a_3b_1)c_1+(a_2b_3-a_3b_2)c_2=0
\]
となる.この式は$pc_1+qc_2=0$の形をしているから,$c_1=q,c_2-p$が解になることがわかる.
これによって,$c_1,c_2$が求まる.そして,式(1)などに代入すれば$c_3$が求まる.
答えを示すと,
\[
    \symbfit{c}=\begin{bmatrix}a_2b_3-a_3b_2 & a_3b_1-a_1b_3 & a_1b_2-a_2b_1\end{bmatrix}^T
\]
\begin{tcolorbox}[colframe=black!80,colback=white,title=外積の定義]
    $\symbfit{a}=\begin{bmatrix}a_1 & a_2 & a_3\end{bmatrix}^T,\symbfit{b}=\begin{bmatrix}b_1 & b_2 & b_3\end{bmatrix}^T$の\textbf{外積}は,ベクトル
    \[
        \symbfit{a}\times \symbfit{b}=\begin{bmatrix}a_2b_3-a_3b_2 & a_3b_1-a_1b_3 & a_1b_2-a_2b_1\end{bmatrix}^T
    \]
\end{tcolorbox}
「これ覚えるの？」と思ったかもしれない.実はこの要素の並び方は不思議なことに\textbf{行列式}になっている!
\\
$\symbfit{i}=\begin{bmatrix}1 & 0 & 0\end{bmatrix}^T,\;\symbfit{j}=\begin{bmatrix}0 & 1 & 0\end{bmatrix}^T,\;\symbfit{k}=\begin{bmatrix}0 & 0 & 1\end{bmatrix}^T$を用いると,
\begin{align*}
    \symbfit{a}\times \symbfit{b}&=(a_2b_3-a_3b_2)\symbfit{i}+(a_3b_1-a_1b_3)\symbfit{j}+(a_1b_2-a_2b_1)\symbfit{k}\\
                                &=\begin{vmatrix}a_2&a_3\\b_2&b_3\end{vmatrix}\symbfit{i}
                                +\begin{vmatrix}a_1&a_3\\b_1&b_3\end{vmatrix}\symbfit{j}
                                +\begin{vmatrix}a_1&a_2\\b_1&b_2\end{vmatrix}\symbfit{k}\\
                                &=\begin{vmatrix}\symbfit{i} & \symbfit{j} & \symbfit{k} \\ a_1 & a_2 & a_3 \\ b_1 & b_2 & b_3\end{vmatrix}\qquad(\text{余因子展開より})
\end{align*}
つまり,この行列式を計算したら,$a\symbfit{i}+b\symbfit{j}+c\symbfit{k}$と値がでる.このときの$a,b,c$が$ \symbfit{a},\symbfit{b}$の共に直交するベクトルの要素となる.
\newpage